\subsection{Methods}
The package \verb+NetworkX+ offers the function \verb+watts_strogatz_graph+ to make a Watts-Strogatz graph. 
Conveniently, this function takes $N$, $k$ and $p$ as arguments. Using this function we created all graphs for this exercise.
Thus, to run our code, the package \verb+NetworkX+ needs to be installed.
We used the suggested values of $N=200$ and $k=6$ for all graphs, but will run several simulations while varying $p$. The different values of the probability will be picked from the following list: $[0.0, 0.001, 0.003, 0.01, 0.03, 0.10, 0.30, 1.0]$.

The graph can also be plotted using the \texttt{draw\_circular()} or \texttt{draw\_networkx()}, which takes in the Watts-Strogatz graph and generates a circular plot or an arbitrary plot. Typically, the graph is shown as a circle and this is how it was presented in figure \ref{fig:WScircleplot}.

\subsubsection{Degrees and degree distribution}

To study how the distribution of the degree $d_i$ of a Watts-Strogatz graph changes by varying $p$ value a pyton script \verb+WS_1_degrees.py+ was written. It can be run using the command \verb+python3 WS_1_degrees.py+. The script creates a Watts-Strogatz graphs with fixed parameters $N=200$ and $k=6$ and a variable parameter $p$ that is varied from $0.0$ to $1.0$. For each $p$ value 1000 separate instances of the graph are created and the degree of each node is studied. Then the results over the 1000 instances are averaged and so are created degree distributions for each $p$ value. Finally histograms are created for each $p$ value and saved as a file \verb+WS_1_degrees.png+. The script utilizes the following Python packages: NetworkX, NumPy and matplotlib.

\input{texfiles/connected_method.tex}

\subsubsection{Diameter and average path length}

The same Python package \texttt{NetworkX} has many useful functions. Once the Watts-Strogatz graph has been generated using the \texttt{nx.watts\_strogatz\_graph} function, the following functions are then used: \texttt{nx.diameter()} and \texttt{nx.average\_shortest\_path\_length()}. As the names of the functions indicate, the \texttt{nx.diameter()} returns the diameter (as defined in equation \ref{eq:WSdiameter}) and the latter function computes the average path length (as defined in equation \ref{eq:WSavgshortestpath}). In the case of an disconnected graph, the diameter and average path length would be computed for the largest connected component of the graph.

The computations are repeated 1000 times for each value of $p$, with a new seed each times. For every iteration, a new graph is generated and the relevant values are saved in a \texttt{NumPy}-array. This will then be used to produce the results visible in section \ref{sec:WS_Diam&ASPL_result}.


\subsubsection{Clustering Coefficient}
The package \verb+NetworkX+ has the function \verb+clustering+, with which the clustering of a graph can be calculated. 
The average shortest path length is calculated using the function \verb+nx.average_shortest_path_length+.
In our code we iterate over the different values of $p$ ($p$ = [0, 0.001, 0.003, 0.01, 0.03, 0.1, 0.3, 1]) as well as 1000 different seeds (1-1000). 
This way the results of our calculations are reproducible. 

The function \verb+networkx.clustering+ outputs a dictionary with the local clustering coefficients of each node. 
For each graph, we take the average of these local clustering coefficients and append them to a list for the current p. 
After iterating over all 1000 seeds, the average for the current p-value is taken and saved to another list.
Additionally, the standard deviation of the clustering coefficients is calculated using numpy and saved to yet another list. 
Then we continue with the next value of p. 

After iterating over all p values, the averaged values and their standard deviation are plotted using \verb+plt.errorbar+. 
All values are normalised by the value at $p=0$.
The whole code is timed, as we were afraid it would take too long to execute it.

To reproduce Figure \ref{fig:cluster_coeff} you need to run \verb+clustering_coeff.py+.
The required packages are \verb+NetworkX+, \verb+numpy+, \verb+matplotlib+ and \verb+time+.
